\section{The two-cycle parameter and escape}\label{sec:witness}

\begin{lemma}\label{lem:witness}
Suppose $a$ is transcendental over $k$ and $f(b)=f(a)+1$. Set
\begin{equation}\label{eq:parameter}
 \lambda_*=1-a-f(a).
\end{equation}
Then $a$ has least period two under $F_{\lambda_*}$. At every
extension $v$ of the pole of $a$ in $k(a)$,
\begin{equation}\label{eq:separation}
 \hhat_{v,\lambda_*}(a)=0,\qquad
 \hhat_{v,\lambda_*}(b)=\frac{\log|a|_v}{36}>0 .
\end{equation}
\end{lemma}

\begin{proof}
Direct expansion in characteristic three gives
\begin{equation}\label{eq:translations}
 \begin{aligned}
 f(-X)&=f(X),\\
 f(X+1)&=f(X)+X-1,\\
 f(X-1)&=f(X)-X-1.
 \end{aligned}
\end{equation}
Consequently the two orbit segments are
\begin{equation}\label{eq:orbits}
 \begin{aligned}
 a&\longmapsto 1-a\longmapsto a,\\
 b&\longmapsto -a-1\longmapsto0\longmapsto\lambda_* .
 \end{aligned}
\end{equation}
For completeness, the nontrivial second images follow from
\[
 \begin{aligned}
 F_{\lambda_*}(1-a)
  &=f(a-1)+1-a-f(a)=-2a=a,\\
 F_{\lambda_*}(-a-1)
  &=f(a+1)+1-a-f(a)=0.
 \end{aligned}
\]
The first image of $b$ uses $f(b)=f(a)+1$ and $2=-1$.
The two points $a$ and $1-a$ are distinct, since their equality
would force $a=-1\in k$. Thus the first orbit has least period two.

Write $C=|a|_v>1$. Since constants are units, the sixth-power
term strictly dominates the other terms at $a$, giving
\begin{equation}\label{eq:radius}
 |f(a)|_v=C^6,\qquad |\lambda_*|_v=C^6.
\end{equation}
If $|z|_v>C$, then $|z|_v^6$ is strictly larger than both
$|z|_v^4$ and $|\lambda_*|_v$. The ultrametric inequality therefore
gives the exact identity
\begin{equation}\label{eq:escape}
 |F_{\lambda_*}(z)|_v=|z|_v^6\qquad(|z|_v>C).
\end{equation}
The second orbit in \eqref{eq:orbits} reaches $\lambda_*$ after
three iterations; by \eqref{eq:radius}, this point is outside that
radius. Induction using \eqref{eq:escape} yields
\[
 |F_{\lambda_*}^n(b)|_v=C^{\,6^{n-2}}\qquad(n\ge3).
\]
Substitution in \eqref{eq:height} gives
\[
 \hhat_{v,\lambda_*}(b)
   =\lim_{n\to\infty}6^{-n}6^{n-2}\log C
   =\frac{\log C}{36}.
\]
The finite first orbit has height zero, proving the lemma.
\end{proof}

\begin{remark}
The orbit identity is uniform in the pair: after the single relation
$f(b)=f(a)+1$ is imposed, $\lambda_*$ depends only on $a$.
The nonzero height does not come from a finite orbit sample; the
strict escape inequality proves it for all subsequent iterations.
The negative difference requires no separate construction: exchange
the roles of $a$ and $b$.
\end{remark}
